\documentclass[10pt,a4paper]{amsart}
\usepackage[margin=19mm]{geometry}
\usepackage{amsmath,amssymb,mathtools}
\usepackage[scr=rsfs,cal=euler]{mathalfa}
\usepackage{xfrac}
\usepackage[unicode,hidelinks,bookmarks=false]{hyperref}
\newcommand{\ie}{i.e.\ }
\newcommand{\cf}{cf.\ }
\newcommand{\resp}{respectively\ }
\newcommand{\Subsection}{\subsection}
\providecommand{\nobreakdash}{\nobreak\hbox{-}}
\setlength{\emergencystretch}{2em}
\multlinegap0pt
\usepackage{amssymb}
\usepackage{algorithm}\usepackage{algpseudocode}
\usepackage{enumerate}
\usepackage{tikz}
\newtheorem{claim}{Claim}
\DeclareMathOperator{\Fix}{Fix}
\DeclareMathOperator{\Int}{int}
\DeclareMathOperator{\Lip}{Lip}
\def\R{\mathbb{R}}
\def\Rk{\mathbb{R}^k}
\def\rr{\mathbf r}
\def\ss{\mathbf s}
\DeclareMathOperator{\supp}{supp}
\def\tt{\mathbf t}
\title{{ \Large\bf{On the representation of measurable and continuous dynamical systems by Lipschitz functions}}}
\author{Yonatan Gutman, Qiang Huo}
\date{Source: arXiv:2505.14653v2 (CC BY 4.0)}
\begin{document}
\maketitle

\noindent\emph{Source and licence.} { \Large\bf{On the representation of measurable and continuous dynamical systems by Lipschitz functions}}. Version arXiv:2505.14653v2 (CC BY 4.0), distributed by arXiv under the Creative Commons Attribution 4.0 International licence, http://creativecommons.org/licenses/by/4.0/, as recorded in the arXiv licence metadata for this exact version and shown on the abstract page https://arxiv.org/abs/2505.14653v2. Full licence notice: \texttt{licenses/arxiv-2505.14653-CC-BY-4.0.txt}.\par\smallskip

\noindent\textit{Editorial scope.} Counted unit: the article's claim perturb (source0.tex lines 739--741). The article's complete original proof of it is delivered with it (source0.tex lines 742--774, one \texttt{proof} environment of 33 lines). No mathematics is written, repaired or re-derived: the statement and the proof are the article's own lines, in the article's own notation, and no retained line is substituted or re-flowed.  No further source-local context is needed: every cross-reference of every retained line resolves inside the two delivered spans.  Nothing was omitted from any retained span, so the package carries no omission record.  No retained line contains a citation, so the wrapper carries no bibliography and no bibliography entry.  No retained line contains a figure environment, an image-inclusion command or drawing code, so no image is delivered and the package declares no asset dependency.  Editorial formatting only, in addition: an AMS \texttt{amsart} preamble that restates the article's own declarations and macros used by the retained lines, namely the environments \texttt{claim} on separate counters, together with the standard packages those lines need.  The retained environments print in this document as Claim 1 in order of appearance, whereas the article prints the same items with the numbering of its own full document.  The complete native source of the article as distributed in the arXiv e-print for this exact version is bundled unchanged as \texttt{sources/178452/source0.tex}.
    \begin{claim}\label{equi and perturb}
It holds that  $g_1\in C_{\Rk,h_0}(X;\Lip_{1}(\Rk))$ and $\max\limits_{x\in X}\max\limits_{\mathbf t\in\Rk}|g_1(x)(\tt)-f(x)(\tt)|\leq 4\delta$. 
    \end{claim}

    \begin{proof}
    It is not hard to see $g_1$ is continuous in the variable $x$. Indeed the key observation is that if $x_i\rightarrow x$ and $\ss \in H(x)$ such that $\ss x \in \supp(q)\subset\Int([-a,a]^k\cdot S)$, then one may find $\ss_i \in \Rk$ such that $\ss_i\rightarrow \ss$ and $\ss_i \in H(x_i)$.  Note that $H(x)=\emptyset$ if $x\in\Fix(X,\Rk)$. Thus, $g_1|_{\Fix(X,\Rk)}=f_1|_{\Fix(X,\Rk)}=h_0$. 
    \noindent
       Let $x\in X$ and $\rr\in \Rk$, if $\mathbf s\in H(\rr x)$, then for $\mathbf t\in[s_1,s_1+a]\times \cdots \times [s_k,s_k+a]$ it holds
    \begin{equation}\label{LHS}
        g_1(\rr x)(\tt)=f_1(\rr x)(\tt)+q((\ss+\rr)x)\cdot u((\ss+\rr)x)(\mathbf {t-s}).
    \end{equation}
    Note that $\mathbf s\in H(\rr x)$ implies that $(\ss+\rr)x\in S$ and therefore $\ss+\rr\in H(x)$ and $\tt+\rr\in [s_1+r_1,s_1+r_1+a]\times \cdots \times [s_k+r_k,s_k+r_k+a]$. Thus
    \begin{equation*}
        g_1(x)(\mathbf {t+r})=f_1(x)(\mathbf {t+r})+q((\ss+\rr)x)\cdot u((\ss+\rr)x)((\mathbf {t+r})-(\mathbf {s+r}))~\text{for}~\mathbf t\in[s_1,s_1+a]\times \cdots \times [s_k,s_k+a],
    \end{equation*}
    which clearly equals to the RHS of Equation \eqref{LHS} as $f_1$ is $\Rk$-equivariant.
    Now assume $\tt\in\Rk\setminus\bigcup\limits_{\mathbf s\in H(\rr x)}[s_1,s_1+a]\times \cdots \times [s_k,s_k+a]$. Thus by definition $g_1(\rr x)(\tt)=f_1(\rr x)(\tt)$. As $f_1$ is $\Rk$-equivariant we have  $g_1(\rr x)(\tt)=f_1(x)(\rr+\tt)$. We claim $g_1(x)(\rr+\tt)=f_1(x)(\rr+\tt)$. Indeed it is enough to show $\rr+\tt\in \Rk\setminus\bigcup\limits_{\mathbf s\in H( x)}[s_1,s_1+a]\times \cdots \times [s_k,s_k+a]$ which follows easily from the condition on $\tt$.  Combining the two cases, we conclude that $g_1$ is $\Rk$-equivariant. Furthermore for all $x\in X$ and $\ss\in H(x)$ and $\tt\in[s_1,s_1+a]\times \cdots \times [s_k,s_k+a]$ it holds
   \begin{equation}
        \begin{split}
    g_1(x)(\tt)&=f_1(x)(\tt)+q(\ss x)\cdot u(\ss x)(\mathbf {t-s})=f_1(x)(\tt)+q(\ss x)\cdot (g(\ss x)(\mathbf {t-s})-f_1(\ss x)(\mathbf {t-s}))\\
    &=(1-q(\ss x))f_1(x)(\tt)+q(\ss x)\cdot g(\ss x)(\mathbf {t-s})\in [0,1].
    \end{split}
     \end{equation}
     We conclude that for all $x\in X$, $g_1(x)(\R^k)\subset [0,1]$.
          Finally observe that for all $x\in X$ and $\tt\in\Rk$
    \begin{equation}\label{eq:g_1-f_1}
        |g_1(x)(\tt)-f_1(x)(\tt)|\leq \|q_{|S}\|_{\infty} \|(g-f_1)_{|X\times [0,a]^k}\|_{\infty}\leq 1\cdot 2\delta
    \end{equation}
      Therefore for all $x\in X$ and $\mathbf t\in\Rk$,
    \begin{equation}\label{eq:g_1-f}
        \begin{split}
            |g_1(x)(\tt)-f(x)(\tt)|&\leq |g_1(x)(\tt)-f_1(x)(\tt)|+|f_1(x)(\tt)-f(x)(\tt)|\\
            &=|g_1(x)(\tt)-f_1(x)(\tt)|+|(1-\delta)f(x)(\tt)+\delta f_0(x)(\tt)-f(x)(\tt)|\\
            &\stackrel{Eq. \eqref{eq:g_1-f_1}}{\leq}2\delta+\delta|f_0(x)(\tt)-f(x)(\tt)|\leq 4\delta.
        \end{split}
    \end{equation}
    \end{proof}

\end{document}
